\section{Methods}
\subsection{Equation of State}
The eCPA can be expressed in terms of the residual Helmholtz free energy as a sum of free energy contributions \cite{Wu1998, Maribo2015}:
\begin{equation}
        A^r = A^\text{phys} + A^\text{assoc} + A^\text{DH} + A^\text{Born},
        \label{Eq:A_eCPA}
\end{equation}
where $A^\text{phys}$, $A^\text{assoc}$, $A^\text{DH}$, and $A^\text{Born}$ represent the free energy contribution from physical interactions, association, Debye-H\"uckel equation and Born term, respectively. When there are no electrolytes in the mixture, the CPA Eos ($A^r = A^\text{phys} + A^\text{assoc}$) is recovered.

Adopting SRK as the cubic EoS, the physical Helmholtz free energy is given by $A^\text{phys} = A^\text{SRK} -A^\text{IG}$, where IG denotes the ideal gas contribution. The SRK Helmholtz free energy \cite{Soave1972} in terms of the molar volume $v$  can be expressed as:
\begin{equation}
        \frac{A^\text{SRK}}{NRT} = - \ln { \left (1-\frac{b}{v}\right )} - \frac{a}{bRT} \ln {\left (1+\frac{b}{v}\right )},
        \label{Eq: Asrk}
\end{equation}
where $N$, $R$, and $T$ are the amount of substance, the universal gas constant and the absolute temperature, respectively, and $a$ and $b$ are the energy and co-volume SRK parameters. In mixtures, the parameters can be obtained from the pure component parameters and their concentration in mole fraction $x_i$ using mixing rules. When the mixture contains only neutral components, the mixture parameters are given by the van der Waals mixing rule:
\begin{equation}
    b = \sum_i b_ix_i,
    \label{Eq:CPA_b}
\end{equation}
\begin{equation}
    a = \sum_j \sum_i x_ix_ja_{ij} \left (1-k_{ij} \right),
    \label{Eq:CPA_a}
\end{equation} 
where $k_{ij}$ is the binary interaction parameter between components $i$ and $j$ ($k_{ij}$ = 0, if $i=j$), and $a_{ij}= \sqrt{a_ia_j}$.

In the presence of electrolytes, the mixture energy parameter is obtained from the Huron-Vidal infinite pressure mixing rule \cite{Maribo2015}:
\begin{equation}
        \frac{a}{b} = \sum_i x_i \frac{a_i}{b_i} - \frac{g^{E\infty}}{\ln{2}}.
        \label{Eq:HuronVidal}
\end{equation}

The excess Gibbs energy $g^{E\infty}$ is obtained from the NRTL model by setting the non-randomness parameter to zero:
\begin{equation}
     g^{E\infty} = \frac{1}{b} \sum_j \sum_i x_ix_jb_j\Delta U_{ji},
    \label{Eq:gEinfity}
\end{equation}
where $\Delta U_{ji}$ is the change in energy between unlike and like interactions ($\Delta U_{ji} = \Delta U_{ij} - \Delta U_{jj}$ ) \cite{Maribo2015}, which is set to zero if $i=j$ and for ion-ion interaction. This parameter can either be made ion-specific or salt-specific, depending on the choice of the parametrization. The first approach allows the modeling of mixed salt solutions, however, it increases the set of parameters to be fitted \cite{Olsen2021}. The interaction between neutral components can still be described by $k_{ij}$ by setting an appropriate  $\Delta U_{ji}$ \cite{Michelsen2004}:
\begin{equation}
     \frac{\Delta U_{ji}}{\ln{2} } = \frac{a_i}{b_i} - \frac{2\sqrt{a_ia_j}}{b_i+b_j} \left (1-k_{ij} \right ),
    \label{Eq:Uij_kij}
\end{equation}
where
$\Delta U_{ij} \neq \Delta U_{ji}$. 
By replacing Eq. \eqref{Eq:Uij_kij} into Eqs. \eqref{Eq:gEinfity} and \eqref{Eq:HuronVidal}, we recover Eq. \eqref{Eq:CPA_a}.

The association Helmholtz free energy $A^\text{assoc}$ is obtained from the Wertheim theory \cite{Wertheim1984}:
\begin{equation}
        \frac{A^\text{assoc}}{NRT} =  \sum_i x_i \sum_{A \in i}  \left ( \ln{\chi_{A_i}} -\frac{\chi_{A_i}}{2} -\frac{1}{2} \right ),
        \label{Eq:Aassoc}
\end{equation}
where $\chi_{A_i}$  is the fraction \textcolor{red}{of molecules $i$ not bonded at the associating site A}. This fraction can be expressed as a function of the compressibility factor $Z$ or molar density $c$ and the association strength between sites ($\Delta_{A_iB_j}$ or $\delta_{A_iB_j}$):
\begin{equation}
    \chi_{A_i} = \frac{Z}{Z+\sum\limits_jx_j \sum\limits_{B \in j} \chi_{B_j}\Delta_{A_iB_j}} = \frac{1}{1+c \sum\limits_j x_j \sum\limits_{B \in j} \chi_{B_j}\delta_{A_iB_j}},
    \label{Eq:CPA_chi}
\end{equation}
\begin{equation}
    \Delta_{A_iB_j} = g_\eta \frac{\kappa_{A_iB_j} P}{RT}\left[ \text{exp}\left( \frac{\varepsilon_{A_iB_j}}{k_BT}\right) -1 \right] = \frac{P}{RT} \delta_{A_iB_j}
    \label{Eq:CPA_delta},
\end{equation}
where \textcolor{red}{$P$ stands for the system pressure,} $\kappa_{A_iB_j}$ and  $\varepsilon_{A_iB_j}$ are association volume and association energy parameters (the volume parameter can also be expressed by $\kappa_{A_iB_j} = b_{ij}\beta_{A_iB_j}$ \cite{Olsen2021}), $g_\eta$ is the radial distribution function, which can be expressed in terms of the packing factor $\eta$: 
\begin{equation}
    g_\eta = \frac{1}{\left(1-1.9\eta\right) },
    \label{CPA_g}
\end{equation}
\begin{equation}
    \eta = \frac{b}{4v}.
    \label{CPA_eta}
\end{equation}

We adopt a 4C association scheme for water molecules, i.e. two proton donors and two proton acceptor sites per molecule. Assuming no \ce{CO2} self-association, but with \ce{CO2}-\ce{H2O} cross-association, \ce{CO2} is modeled as a ``pseudo-associating'' species, with the same 4C association scheme as water and the cross-association strength described as a fraction of the water self-association ($\delta_{A_wB_c} = S_{cw}\delta_{A_wB_w}$) \cite{Li2009}. Cross-association is assumed to be symmetric between two sites of water and \ce{CO2} ($\delta_{A_wB_c} = \delta_{B_wA_c}$). In that case, we can omit the site subscripts, and the non-bonded fraction of water and \ce{CO2} molecules are given, respectively, by:
%\begin{equation}
%    \chi_w = \frac{Z}{Z+2x_w\chi_w\Delta_w+2x_c\chi_cS_{cw}\Delta_w}, 
%    \label{chiw_withCO2}
%\end{equation}
%\begin{equation}
%    \chi_c = \frac{Z}{Z+2x_w\chi_wS_{cw}\Delta_w}.
%    \label{chic}
%\end{equation}

\textcolor{red}{
\begin{equation}
     \chi_w = \frac{1}{1+2cx_w\chi_w\delta_w+2cx_c\chi_cS_{cw}\delta_w}, 
    \label{chiw_withCO2}
\end{equation}
\begin{equation}
    \chi_c = \frac{1}{1+2cx_w\chi_wS_{cw}\delta_w}.
    \label{chic}
\end{equation}
}

The DH Helmholtz free energy is given by \cite{Debye1923}: 
\begin{equation}
        \frac{A^\text{DH}}{NRT} = -\frac{\sum_j x_j z_j^2 X_j}{4\pi N_A c \sum_j x_j z_j^2}, 
        \label{Eq:ADH}
\end{equation}
where $N_A$ is the Avogadro constant and $z_j$ is the valence of ion $j$. The function $X_j$ depends on the hard-sphere diameter $\sigma_j$ of ion $j$ and the reciprocal Debye length $\kappa$, which is a function of the dielectric constant of the solvent $\epsilon_r$, the vacuum permittivity $\epsilon_0$, the Boltzmann constant $k_B$ and the elementary charge $e$:
\begin{equation}
        X_i = \frac{1}{\sigma_i^3} \left [ \ln { \left (1 + \kappa \sigma_i \right )} - \kappa \sigma_i + \frac{\left ( \kappa \sigma_i \right)^2}{2} \right ],
        \label{Eq:ADH_Xi}
\end{equation}
\begin{equation}
        \kappa = \sqrt{\frac{e^2c\sum_jx_jz_j^2}{k_BT\epsilon_0\epsilon_r}}.
        \label{Eq:ADH_debye}
\end{equation}

The Born Helmholtz free energy is given by \cite{Michelsen2004}: 
\begin{equation}
        \frac{A^\text{Born}}{NRT} = \frac{N_A e^2}{8\pi \epsilon_0RT } \left ( \frac{1}{\epsilon_r} - 1\right ) \sum_j \frac{x_jz_j^2}{R_{B,j}},
        \label{ABorn}
\end{equation}
where $R_{B,j}$ is the radius of the Born cavity, which is different from the hard-sphere diameter because it accounts for the ion solvation sphere.

To describe the DH and Born free energy, a relative permittivity model is required. By coupling Onsager and Kirkwood theories, the static permittivity in mixtures can be obtained from the infinite frequency permittivity $\epsilon_\infty$, the Kirkwood g-factor $g_i$ and the vacuum dipole moment $\mu_{0i}$ of component $i$ \cite{Davies1972}:
\begin{equation}
        \frac{\left (2\epsilon_r + \epsilon_\infty \right ) \left ( \epsilon_r - \epsilon_\infty \right )}{\epsilon_r\left ( \epsilon_\infty + 2 \right )^2} = \frac{N_A}{9\epsilon_0k_BTv}\sum_jx_jg_j\mu_{0j}^2.
        \label{Eq:eps}
\end{equation} 

The infinite frequency permittivity is obtained from Clausius-Mossotti equation from the molecular polarizability $\alpha_{0i}$ of each molecule \cite{Scaife1989}:
\begin{equation}
     \frac{\epsilon_\infty-1}{\epsilon_\infty+2}=\frac{N_A}{3\epsilon_0v}\sum_jx_j\alpha_{0j}.
        \label{Eq:eps_infty}
\end{equation} 

The Kirkwood g-factor depends on the coordination number and the statistical distribution of dipole moment angles. By adopting a primitive model, \citeauthor{Maribo2013a} \cite{Maribo2013a} have proposed a correlation between the Kirkwood g-factor and the Wertheim association theory. The g-factor of polar components can be expressed as: 
\begin{equation}
        \color{red}
        g_i = 1 + \sum_j \frac{z_{ij}\mathcal{P}_{ij}\cos{\gamma_{ij}}}{\mathcal{P}_i \cos{\theta_{ij}+1}}\left ( \frac{\mu_{0j}}{\mu_{0i}} \right )
        \label{Eq:Kirkwood}
\end{equation} 
where $z_{ij}$, \textcolor{red}{ $\mathcal{P}_{ij}$}, $\gamma_{ij}$ and $\theta_{ij}$ are, respectively, the coordination number, the probability of association, the dipole moments angle and the shells rotation angle between molecules $i$ and $j$, and \textcolor{red}{ $\mathcal{P}_{i}$} is the probability the molecule $i$ is associating with any molecules \textcolor{red}{ ($\mathcal{P}_{i} = \sum\limits_j \mathcal{P}_{ij}$)}. The probability of association between molecules depends on the nonbonded fraction of sites from the Wertheim theory:
\textcolor{red}{
\begin{equation}
        \mathcal{P}_{ij} = \mathcal{N}_{A_i} c x_j\chi_{A_i} \chi_{B_j} \delta_{A_iB_j},
        \label{Eq:P_ij}
\end{equation} 
}
where $\color{red}\mathcal{N}_{A_i}$ is the number of sites $A$ in molecule $i$.

To determine density and the phase equilibria of \ce{CO2}-water and \ce{CO2}-brine mixtures from eCPA EoS, we need expressions for the compressibility factor and fugacity coefficients, which are presented in the Supporting Information.

\subsubsection{Parametrization}
The pure component parameters of \ce{CO2} and \ce{H2O} are taken from \citeauthor{Tsivintzelis2011} \cite{Tsivintzelis2011}, which are based on experimental vapor pressure and liquid density. The parameter $a_i$ is made temperature dependent using the critical temperature, \textcolor{red}{$T_{c,i}$, of component $i$}:
\begin{equation}
    \color{red}
    a_i = a_{0i}\left [ 1+ c_{1i} \left ( 1-\sqrt{\frac{T}{T_{c,i}}} \right )\right ]^2.
    \label{Eq:CPA_ai}
\end{equation}

The ion parameters are from \citeauthor{Maribo2015} \cite{Maribo2015}. The authors have used hydration enthalpy, osmotic coefficient, and mean ionic activity coefficient data in the regression. To decrease the number of adjustable parameters for ions, the SRK energy parameter is set to zero, the co-volume is obtained from the hard-sphere diameter ($b_i=2N_A \pi \sigma_i^3/3$),  and the interaction with neutral components is purely electrostatics (no association). To properly represent the brine density, a Peneloux volume translation is required ($v = v^\text{EoS} + x_s \Upsilon_s$, where $\Upsilon_s$ is the Peneloux salt parameter) \cite{Maribo2015}.  The benefit of using such strategy is the improvement in density predictions, without affecting properties such as activity coefficients and phase equilibria \cite{Olsen2021}. The \ce{CO2}-water mixture requires no volume translation. Table S1 of the Supporting Information presents the parameters for the pure components.


The binary interaction parameters are made temperature-dependent to increase the condition range of the model. For \ce{CO2}-\ce{H2O} mixture, the binary interaction, and the cross-association parameters are given, respectively, by:
\textcolor{red}{
\begin{equation}
    k_{cw} = A_k \left ( {\frac{T}{T_{c,c}}} \right )^2 + B_k \left ( {\frac{T}{T_{c,c}}} \right ) + C_k,
    \label{Eq:kij}
\end{equation}
\begin{equation}
    S_{cw} = A_S \left ( {\frac{T}{T_{c,c}}} \right )^2 + B_S \left ( {\frac{T}{T_{c,c}}} \right ) + C_S,
    \label{Eq:Sij}
\end{equation}
where $T_{c,c}$ is the \ce{CO2} critical temperature (304 K) and} $A_k$, $B_k$, $C_k$,  $A_S$, $B_S$, and $C_S$ are adjusted based on vapor-liquid (VLE) and liquid-liquid (LLE) equilibria data from the \ce{CO2}-water mixture. To compensate for the limited data available at high temperatures, we lump the error in the objective function (O.F.) in sets of temperature ranges. 
%Otherwise, the large amount of data at low temperatures would benefit the model only in this range. 
Another strategy to improve the predictions at higher temperatures is to use the \ce{CO2} concentration in both aqueous ($x_c^W$) and \ce{CO2}-rich phase ($x_c^C$):
\begin{equation}
    \text{O.F.} = \min \sum_k \left (\frac{1}{N^W_k} \sum_i \left |\frac{ x_{c,i}^{W,\text{calc}}-x_{c,i}^{W,\text{exp}}}{ x_{c,i}^{W,\text{exp}}} \right | + \frac{1}{N^C_k} \sum_i \left |\frac{ x_{c,i}^{C,\text{calc}}-x_{c,i}^{C,\text{exp}}}{ x_{c,i}^{C,\text{exp}}} \right | \right ),
    \label{Eq:OF1}
\end{equation}
where in the first summation $k$ runs in the subset of temperatures, which upper limits are given by \{298, 323, 353, 393, 423, 473, 538, 573, 598, 623 K\} and $N_k^\alpha$ is the number of experimental data in the temperature subset $k$ in phase $\alpha$. The CPA EoS is adjusted in the temperature and pressure ranges of 283 to 623 K and 5 to 3500 bar. Table \ref{EXPs_CW} shows the range of conditions and the average absolute deviation of each phase for each experimental work, whereas Table \ref{Kij_Sij} shows the optimized parameters for the \ce{CO2}-water mixture.


\begin{table}[]
\caption{Experimental data set of \ce{CO2}-\ce{H2O} phase equilibria applied in the CPA parametrization with the respective temperature and pressure range, number of data points $N_p$, and average absolute deviation AAD of the calculated mutual solubility for aqueous ($W$) and \ce{CO2}-rich ($C$) phases.}
\centering
\begin{tabular}{lcccccc}
\hline
\hline
 \textbf{Reference} & $T$ (K) & $P$ (bar)  & $N_p^w$  & AAD$^W$ (\%)$^a$ & $N_p^C$ & AAD$^C$ (\%)$^b$ \\ 
 \hline
\citeauthor{Guo2014} \cite{Guo2014} & 283 — 573 & 100 — 1200 & 122 & 6.2  & 0 & - \\
\citeauthor{Meyer2015} \cite{Meyer2015} &  283  — 353 & 5  — 50 & 0 & - & 58 & 14.2 \\
\citeauthor{Valtz2004} \cite{Valtz2004} &  283 — 318 &  5 — 80 & 38 & 13.6 & 30 & 19.4 \\
\citeauthor{King1992} \cite{King1992} &  50 — 250 & 288 — 313 & 18 & 3.6 & 26 & 74.0 \\
\citeauthor{Hou2013} \cite{Hou2013} & 298 — 448 & 10 — 175 & 36 & 4.6 & 32 & 10.7 \\
\citeauthor{Briones1987} \cite{Briones1987} & 323 &  70 — 175 & 7 & 9.6 & 7 & 10.7 \\
\citeauthor{Bamberger2000} \cite{Bamberger2000} & 323 — 353 & 40 — 140 & 29 & 4.4 & 29 & 16.4 \\
\citeauthor{Rumpf1994} \cite{Rumpf1994} & 323 & 10 — 60 & 5 & 4.3 & 0 & - \\
\citeauthor{Dsouza1988} \cite{Dsouza1988} & 323 — 348 & 100 — 150 & 4 & 10.2 & 4 & 24.5 \\
 \citeauthor{Dohrn1993} \cite{Dohrn1993} &  323 & 100 — 300 & 3 & 8.3 & 3 & 31.8 \\
\citeauthor{Tödheide1963} \cite{Tödheide1963} & 323 — 623 & 200 — 3500 & 98 & 9.3 & 98 & 34.5 \\
\citeauthor{Sako1991} \cite{Sako1991} & 348 — 423 & 100 — 200 & 7 & 15.1 & 8 & 21.0 \\
\citeauthor{Nighswander1989} \cite{Nighswander1989} & 353 — 473 & 20 — 100 & 19 & 18.2 & 0 & - \\
\citeauthor{Muller1988} \cite{Muller1988} & 373 — 473 & 5 — 80 & 40 & 17.8 & 0 & - \\
\citeauthor{Tong2013} \cite{Tong2013} &  375 & 70 — 275 & 5 & 3.0 & 0 & - \\
\citeauthor{Takenouchi1964} \cite{Takenouchi1964} & 383 — 623 & 100 — 1500 & 108 & 7.0 & 108 & 21.3 \\ 
\hline 
\textbf{Global} & \textbf{283 — 623} & \textbf{5 — 3500} & \textbf{539} & \textbf{8.6} & \textbf{443} & \textbf{23.9} \\ 
 \hline
 \hline
\end{tabular}
%\raggedleft
\flushleft
\footnotesize{$^a$ AAD$^W$ = 1/$N_p^W$ $\sum |x_c^{W,\text{calc}}-x_c^{W,\text{exp}}|$/ $x_c^{W,\text{exp}}$} 

\footnotesize{$^b$ AAD$^C$ = 1/$N_p^C$ $\sum |x_w^{C,\text{calc}}-x_w^{C,\text{exp}}|$/ $x_w^{C,\text{exp}}$}
\label{EXPs_CW}
\end{table}


\begin{table}[]
\caption{Optimized binary ($k_{cw}$) and cross-interaction ($S_{cw}$) parameters for the \ce{CO2}-\ce{H2O} mixture with CPA EoS}
\centering
\begin{tabular}{cccc}
\hline
\hline
 \textbf{Parameter} & $A_l$ & $B_l$ & $C_l$\\ 
 \hline
$k_{cw}$ & -0.492 & 2.101 & -1.571 \\
$S_{cw}$ & 0.192 & -0.173 & -0.009 \\
 \hline
 \hline
\end{tabular}
\label{Kij_Sij}
\end{table}

\textcolor{red}{If ion-specific mixture parameters were chosen, more adjustable variables would be required. Because the mixture in our work has a single salt (NaCl)}, the interaction parameters between neutral and charged species are salt-specific and made temperature-dependent by:
\textcolor{red}{
\begin{equation}
    \Delta U_{is} = \Delta U^\text{ref}_{is} + \alpha_{is}R \left [ \left ( 1-\frac{T}{T^\alpha_{i,s}} \right )^2 -\left ( 1-\frac{T^\text{ref}}{T^\alpha_{i,s}} \right )^2 \right ],
    \label{Eq:DUij}
\end{equation}
}
where $i$ = $w$ or $c$, and $T^\text{ref}$ = 298 K. The interaction parameter between water and NaCl is taken from \citeauthor{Maribo2015} \cite{Maribo2015}; the interaction parameter between \ce{CO2} and NaCl is adjusted accounting for \ce{CO2}-\ce{H2O} cross-association using \ce{CO2} solubility in NaCl brine ($m_c$) experimental data in the regression. The objective function is given by: 
\begin{equation}
    \text{O.F.} = \min  \sum_i \left |\frac{ m_{c,i}^{\text{calc}}-m_{c,i}^{\text{exp}}}{ m_{c,i}^{\text{exp}}} \right |.
    \label{Eq:OF2}
\end{equation}

The eCPA EoS is adjusted in the temperature, pressure and NaCl salinity ranges of 283 to 723 K, 1 to 1400 bar and 0 to 6 mol$\cdot$kg$^{-1}$. Table \ref{EXPs_CB} shows the range of conditions (temperature, pressure and salinity) and the average absolute deviation of \ce{CO2} solubility in brine for each experimental work, whereas Table \ref{Ucs_Uws} shows the optimized parameters for the \ce{CO2}-brine mixture.


\begin{table}[]
\caption{Experimental data set of \ce{CO2}-\ce{NaCl}-\ce{H2O} phase equilibria applied in the eCPA parametrization with the respective temperature, pressure and salinity range, number of data points $N_p$, and average absolute deviation AAD of the calculated \ce{CO2} solubility in brine.}
\centering
\begin{tabular}{lccccc}
\hline
\hline
 \textbf{Reference} & $T$ (K) & $P$ (bar)  & $m_s$ (mol$\cdot$kg$^{-1}$) &$N_p$  & AAD (\%)$^a$\\ 
 \hline
\citeauthor{Guo2016} \cite{Guo2016} & 283 — 473 & 100 — 400 & 1 — 5 & 156  & 4.5 \\
\citeauthor{Liu2021} \cite{Liu2021} & 298 — 373 & 1 — 200 & 0.17 — 0.26 & 81  & 6.1 \\
\citeauthor{Bando2003} \cite{Bando2003} & 303 — 333 & 100 — 200 & 0.17 — 0.52 & 36  & 9.4 \\
\citeauthor{Chabab2019} \cite{Chabab2019} & 303 — 373 & 35 — 350 & 6 & 14  & 7.2 \\
\citeauthor{Kiepe2002} \cite{Kiepe2002} & 313 — 353 & 1 — 100 & 0.52 — 4 & 52  & 5.8 \\
\citeauthor{Rumpf1994} \cite{Rumpf1994} & 313 — 413 & 5 — 100 & 4 & 63  & 3.8 \\
\citeauthor{Liu2011b} \cite{Liu2011b} & 318 & 20 — 160 & 2 & 8  & 3.5 \\
\citeauthor{Hou2013} \cite{Hou2013} & 323 — 423 & 25 — 180 & 2.5 — 4 & 36  & 9.3 \\
\citeauthor{Koschel2006} \cite{Koschel2006} & 323 — 373 & 50 — 200 & 1 — 3 & 14 & 5.9 \\
\citeauthor{Messabeb2016} \cite{Messabeb2016} & 323 — 423 & 50 — 200 & 1 — 6 & 36  & 6.3 \\
\citeauthor{Yan2011} \cite{Yan2011} & 323 — 413 & 50 — 400 & 1 — 5 & 36  & 6,7 \\
\citeauthor{Zhao2015} \cite{Zhao2015} & 323 — 423 & 150 & 1 — 6 & 18  & 4.5 \\
\citeauthor{Nighswander1989} \cite{Nighswander1989} & 353 — 473 & 20 — 100 & 0.17 & 18  & 17.2 \\
\citeauthor{Savary2012} \cite{Savary2012} & 423 & 120 — 340 & 2 & 6  & 6.7 \\
\citeauthor{Takenouchi1965} \cite{Takenouchi1965} & 423 — 723 & 100 — 1400 & 1.09 — 4.27 & 86  & 12.2 \\
\hline 
\textbf{Global} & \textbf{283 — 723} & \textbf{1 — 1400} & \textbf{0.17 — 6} & \textbf{660} & \textbf{6.9} \\ 
 \hline
 \hline
\end{tabular}
\flushleft
\footnotesize{$^a$AAD = 1/$N_p$ $\sum |m_c^{\text{calc}}-m_c^{\text{exp}}|$/ $m_c^{\text{exp}}$} 
\label{EXPs_CB}
\end{table}

\begin{table}[]
\caption{Optimized interaction parameters between water-salt ($\Delta U_{ws}$)\cite{Maribo2015} and \ce{CO2}-salt ($\Delta U_{cs}$) for the \ce{CO2}-\ce{NaCl}-\ce{H2O} mixture with eCPA EoS}
\centering
\begin{tabular}{cccc}
\hline
\hline
 \textbf{Parameter} & $\Delta U^\text{ref}_{is}$ (J.mol$^{-1}$) & $\alpha_{is}$ (K) & \textcolor{red}{$T^\alpha_{i,s}$} (K)\\ 
 \hline
$\Delta U_{ws}$ & -1858 & 1573 & 340 \\
$\Delta U_{cs}$ & 6056 & 692 & 244 \\
 \hline
 \hline
\end{tabular}
\label{Ucs_Uws}
\end{table}



\newpage

\subsection{Molecular Simulations}
Molecular simulations are carried out in the Gibbs ensemble Monte Carlo with fixed temperature and pressure (NPT-GEMC) using the open software Cassandra \cite{cassandra}. A total of 2000 molecules is placed in the two-box system. The \ce{CO2} global composition varies between 15 and 30\% (mole basis) depending on the condition. To avoid energy drift, first, each box is simulated individually in a short (1 million Monte Carlo steps) NPT ensemble close to the expected equilibrium phase composition. Then, 50 million MC steps are performed in the GEMC equilibration stage, in which the maximum displacement allowed is adjusted to match 50\% acceptance rates in MC moves. The production is executed for more 50 million MC steps with fixed maximum displacement width, and information is collected on the density and composition of both phases. The probability of attempting a translation, rotation, particle transfer (swap), and volume exchange move is set at 35, 30, 34.75, and 0.25\%, respectively. In the \ce{CO2}-brine mixture, the probability of transferring ions to the \ce{CO2}-rich phase is set to zero because ions remain in the aqueous phase. 

We investigate the \ce{CO2}-water phase equilibria at $T$ = 548, 573, 598 and 623 K and pressures from 200 to 1200 bar. At least 3 independent simulations are performed for each condition and the statistical errors are determined based on the blocking average of each simulation and the standard deviation of the independent simulations. Close to the critical point, macroscopic fluctuations may occur and the boxes may switch their phase representation, which requires more independent simulations. We keep track of the energy, density, and composition of boxes to guarantee they are equilibrated and represent the right phase. The \ce{CO2}-brine mixture is investigated at  $T$ = 573 and 623 K, and $P$ = 200, 300, 400, 500 and 600 bar. Although there is no salt in the \ce{CO2}-rich phase, water molecules are distributed in both boxes, in such a way the salinity is not kept constant beforehand \cite{Vorholz2004}.  

\ce{CO2} is represented by a semiflexible (rigid bonds with an harmonic potential in the angle) EPM2 \cite{Harris1995} and water by a rigid SPCE \cite{Berendsen1987}. This combination has not been investigated at high temperatures in the past. For the \ce{Na+} and \ce{Cl-} ions, a classical integer charge force field is applied \cite{Smith1994}. The force field parameters and partial charges are shown in Table S2 of the Supporting Information. The non-bonded interactions between particles account for the van der Waals and electrostatics interactions, which are represented by the Lennard-Jones (LJ) and Coulombic potential, respectively. We investigate the sensibility of the simulations on the unlike LJ parameters between \ce{CO2} and \ce{H2O}: they are either obtained by applying Lorentz-Berthelot (LB) combining rules or using optimized (OPT) parameters by \citeauthor{Vlcek2011} \cite{Vlcek2011} for the \ce{CO2}-\ce{H2O} phase equilibria at temperatures lower than 350 K \textcolor{red}{(Table S3 of the Supporting Information)}. The LJ potential is cut at a distance of 1.2 nm and analytical tail corrections are applied to the energy and pressure. The electrostatics are handled with Ewald summation with an accuracy set at 10$^{-5}$.

